\documentclass[11pt,a4paper]{scrartcl} 

% Kalbos nustatymai (lietuvių kalbai)
\usepackage[utf8]{inputenc}
\usepackage[T1]{fontenc}
\usepackage{tasks}
\usepackage{tikz}
\usepackage{lmodern} % Įkeliame modernų šriftą, kuris palaiko visas reikalingas formas
\usepackage[lithuanian]{babel}

% Įkeliame jūsų .sty failą (įsitikinkite, kad Overleaf projekte šalia yra failas manostilius.sty)
\usepackage{../manostilius_lt}

% Projekto informacija (būtina preambulėje dėl jūsų stiliaus specifikos)
\title{Algebriniai reiškiniai}
\author{Andrius Berniukevičius} 

\begin{document}

\maketitle


\section{Algebrinių reiškinių tipai}

Šiame dokumente pateikiami jūsų \texttt{.sty} faile apibrėžtų aplinkų ir makrokomandų pavyzdžiai. Pradėkime nuo skaičių aibių. Mes žinome, kad natūralieji skaičiai $\NN$, sveikieji skaičiai $\ZZ$, racionalieji skaičiai $\QQ$ ir realieji skaičiai $\RR$ tenkina sąryšį:
\[ \NN \subset \ZZ \subset \QQ \subset \RR \subset \CC \]
\begin{definition}
    \vocab{Algebrinis reiškinys} – tai matematinis užrašas, sudarytas iš skaičių ir kintamųjų (raidžių), sujungtų aritmetinių veiksmų (sudėties, atimties, daugybos, dalybos, kėlimo laipsniu ar šaknies traukimo) ženklais. Reiškiniuose taip pat gali būti naudojami skliaustai, nurodantys veiksmų atlikimo tvarką.
\end{definition}

\begin{example}
Štai atskiri algebrinių reiškinių pavyzdžiai, iliustruojantys kiekvieną pagrindinį matematinį veiksmą:
\begin{itemize}
    \item \textbf{Sudėtis ir atimtis:} 
    \[ 3x + 5y - 7 \]
    
    \item \textbf{Daugyba:} 
    \[ 4a(a - 2b) \]
    
    \item \textbf{Dalyba (trupmena):} 
    \[ \frac{2x - 1}{x + 4} \]
    
    \item \textbf{Kėlimas laipsniu:} 
    \[ 5x^4 - 2x^2 + 8 \]
    
    \item \textbf{Šaknies traukimas:} 
    \[ \sqrt{x^2 + y^2} \]
\end{itemize}
\end{example}

Algebriniai reiškiniai skirstomi į \textbf{racionalius} ir \textbf{irracionalius} reiškinius pagal tai, ar su bent vienu kintamuoju yra atliekamas šaknies traukimo veiksmas.
\begin{definition}
    \vocab{Racionalus reiškinys} - tai algebrinis reiškinys, kuriame su kintamaisiais (raidėmis) atliekami tik šie veiksmai: sudėtis, atimtis, daugyba, dalyba ir kėlimas sveikuoju laipsniu. Juose nėra šaknies traukimo iš kintamojo veiksmo.
\end{definition}
\begin{example}
Racionaliųjų reiškinių pavyzdžiai:
    \begin{itemize}
        \item $3x ^{2} -5$
        \item $\frac{a+b}{c} $
        \item $(x-2)(x+3)$ 
    \end{itemize}
\end{example}

\begin{definition}
    \vocab{Irracionalus reiškinys} - tai algebrinis reiškinys, kuriame bent vienas kintamasis (raidė) yra po šaknies ženklu arba yra keliamas trupmeniniu laipsniu.
\end{definition}
\begin{example}
Racionaliųjų reiškinių pavyzdžiai:
    \begin{itemize}
        \item $5\sqrt{a}+5b $
        \item $ \frac{4}{\sqrt[3]{x+y} }  $
        \item $x ^{\frac{1}{3}} +4x-5   $ 
    \end{itemize}
\end{example}
\begin{remark}
    Reiškinys $\sqrt{5}x-3y$ nėra irracionalus, nes po šaknimi yra skaičius, o ne kintamasis.
\end{remark}

\textbf{Racionalieji reiškiniai} yra skirstomi \textbf{sveikus} ir \textbf{trupmeninius} reiškinius pagal tai, ar juose yra dalyba iš kintamojo.


\begin{definition}
    \vocab{Sveikas reiškinys} - tai racionalus reiškinys, kuris neturi dalybos iš kintamojo.
\end{definition}

\begin{example}
Sveikų reiškinių pavyzdžiai:
    \begin{itemize}
        \item $7x ^{3} -4x+1$
        \item $(a+2b) ^{3} $
        \item $\frac{x+3}{7} $ 
    \end{itemize}
\end{example}

\begin{definition}
    \vocab{Trupmeninis reiškinys} - tai racionalus reiškinys, kuriame yra dalyba iš kintamojo
\end{definition}

\begin{example}
Trupmeninių reiškinių pavyzdžiai:
    \begin{itemize}
        \item $\frac{5}{x}+x-8$
        \item $\frac{a+b}{a-b}  $
        \item $\frac{5}{x-3}+ \frac{1}{x}  $ 
    \end{itemize}
\end{example}
\begin{remark}
    Nors reiškinyje $4x^{-2} $ dalybos tiesiogiai nematome, tai nėra racionalus reiškinys, nes:
    $$4x ^{-2} = 4 \cdot \frac{1}{x ^{2} }  = \frac{4}{x ^{2} }   $$
    Kadangi vardiklyje yra kintamasis, vadinasi, tai yra trupmeninis reiškinys.
\end{remark}

Iš sveikųjų reiškinių išskirkime \textbf{vienanarius} ir \textbf{daugianarius}.
\begin{definition}
    \vocab{Vienanaris} - tai toks sveikasis reiškinys, kuriame su kintamaisiais ir skaičiais neatliekami jokie sudėties $(+)$  arba atimties $(-)$ veiksmai. Kitaip tariant, tai yra sveikasis reiškinys, kurį sudaro tik skaičių, kintamųjų ir jų natūraliųjų laipsnių sandauga.
\end{definition}
\begin{example}
Vienanarių pavyzdžiai:
\begin{tasks}[style = itemize](3)
    \task $5a$
    \task $-3a ^{2} b ^{4} $
    \task $xyz$
    \task $\frac{1}{2}x ^{2} yz $ 
    \task $\frac{2a}{5} $
    \task $8$ 
\end{tasks}

\end{example}
\begin{remark}
 Trupmena $\frac{2 a}{5}$ laikoma vienanariu tik todėl, kad ją galima užrašyti kaip sandaugą: $\frac{2}{5} \cdot a$, o toks reiškinys pagal apibrėžimą yra vienanaris. 
\end{remark}

\begin{remark}
Skaičius $8$ yra laikomas vienanariu, nes jį galima užrašyti kaip sandaugą $8 \cdot 1$, o toks reiškinys pagal apibrėžimą yra vienanaris.
\end{remark}

\begin{definition}
    \vocab{Daugianaris} - sveikasis reiškinys, sudarytas iš vienanarių algebrinės sumos.
\end{definition}

\begin{example}
    \begin{tasks}[style = itemize](2)
        \task $x ^{2} +3x-4$
        \task $ab+bc+cd$
        \task $x ^{3} +xy-y ^{2}+1 $
        \task $x+8$ 
    \end{tasks}
\end{example}


Visų algebrinių reiškinių klasifikacijos medis:

\begin{center}
\begin{tikzpicture}[
  % Atstumai tarp lygių ir šakų
  level distance=2.5cm,          % Šiek tiek padidintas gylis, kad tekstas turėtų erdvės
  level 1/.style={sibling distance=8.0cm},
  level 2/.style={sibling distance=4.5cm},
  level 3/.style={sibling distance=4.5cm}, % Padidintas atstumas tarp vienanarių ir daugianarių
  % Bendras mazgų stilius
  every node/.style={draw, rectangle, rounded corners, text width=3.6cm, align=center, thick},
  % Specifiniai stiliai skirtingiems lygiams (su spalvomis)
  root/.style={fill=blue!15, text width=4.5cm, font=\bfseries},
  level1/.style={fill=green!10, font=\bfseries},
  level2/.style={fill=yellow!15},
  level3/.style={fill=orange!15, text width=3.8cm} % Praplatinti langeliai papildomam tekstui
]

  % Medžio braižymas
  \node[root] {Algebriniai reiškiniai}
    child {node[level1] {Racionalieji reiškiniai \\ \small(be kintamojo po šaknimi)}
      child {node[level2] {Sveikieji reiškiniai \\ \small(be kintamojo vardiklyje)}
        child {node[level3] {Vienanariai \\ \small(tik skaičių ir raidžių sandauga, jokių $+/-$)}}
        child {node[level3] {Daugianariai \\ \small(vienanarių algebrinė suma}}
      }
      child {node[level2] {Trupmeniniai reiškiniai \\ \small(su kintamuoju vardiklyje)}}
    }
    child {node[level1] {Iracionalieji reiškiniai \\ \small(kintamasis yra po šaknimi)}};

\end{tikzpicture}
\end{center}


\section{Uždaviniai}
\begin{problem}
Prie kiekvieno žemiau esančio reiškinio parašykite, kokio tipo tai reiškinys: racionalus, iracionalus, sveikas, trupmeninis, vienanaris ar daugianaris. Atsiminkite, kad reiškinys gali turėti daugiau nei vieną tipą.

\begin{tasks}[ label={(\arabic*)},label-offset=0.9em](4) % Nustatome skaičiavimą ir formatą (1)
    \task $7x^3 -x- 15$
    \task $\sqrt{ab} + 5$
    \task $\frac{7x}{3}$
    \task $\frac{5}{y ^{3} }$
    %
    \task $-6a^3bc ^{2} $
    \task $x^2 + 2xy + y^2$
    \task $\sqrt{7}x - 2x ^{4} $
    \task $5a^{-3}$
    %
    \task $12$
    \task $\frac{x+3}{x-3}$
    \task $y^{\frac{2}{3}} + 5y$
    \task $\frac{x-1}{3} + 4x$
    %
    \task $abcd$
    \task $\frac{1}{\sqrt[3]{x}+2}$
    \task $(x-2)^3$
    \task $-\frac{3}{5}m^2n$
    %
    \task $x + 8$
    \task $(x+3)(x-2)$
    \task $\frac{4}{x^2}$
    \task $a^2 - b^2$
\end{tasks}
\end{problem}





\end{document}